\documentclass[10pt,a4paper]{amsart}
\usepackage{amsmath,amssymb,amsthm,graphicx}
% Retained preamble for the Stacks Project chapter snapshot in this directory.
% The chapter's own preamble.tex is a copy of the shared Stacks Project preamble; the block
% below is sliced MECHANICALLY, line for line, from it. Nothing was added, removed or
% redefined; the wrapper class, the xr-hyper/hyperref wiring and the externaldocument list
% are omitted because the assembler supplies the document class.
%
% lines taken from preamble.tex: 3, 5, 9, 11, 12, 156, 157, 158, ... (51 lines total)
%
% package coverage restored from the chapter's own preamble (2026-10-13): xy, xy-2cell
% package coverage restored from the chapter's own preamble (2026-10-14): verbatim
%   the chapter's preamble.tex loads \usepackage{verbatim} on line 8 and then defines
%   reference/slogan/history as \comment environments; without the package \comment is
%   undefined and any card whose counted spans use one of the three environments fails
%   with "Undefined control sequence" (L44184, lemma-representable-over-space).

\usepackage{verbatim}
\usepackage[all]{xy}
\xyoption{2cell}
\UseAllTwocells

\newenvironment{reference}{\comment}{\endcomment}
\newenvironment{slogan}{\comment}{\endcomment}
\newenvironment{history}{\comment}{\endcomment}
\theoremstyle{plain}
\newtheorem{theorem}[subsection]{Theorem}
\newtheorem{proposition}[subsection]{Proposition}
\newtheorem{lemma}[subsection]{Lemma}
\theoremstyle{definition}
\newtheorem{definition}[subsection]{Definition}
\newtheorem{example}[subsection]{Example}
\newtheorem{exercise}[subsection]{Exercise}
\newtheorem{situation}[subsection]{Situation}
\theoremstyle{remark}
\newtheorem{remark}[subsection]{Remark}
\newtheorem{remarks}[subsection]{Remarks}
\numberwithin{equation}{subsection}
\def\lim{\mathop{\mathrm{lim}}\nolimits}
\def\colim{\mathop{\mathrm{colim}}\nolimits}
\def\Spec{\mathop{\mathrm{Spec}}}
\def\Hom{\mathop{\mathrm{Hom}}\nolimits}
\def\Ext{\mathop{\mathrm{Ext}}\nolimits}
\def\SheafHom{\mathop{\mathcal{H}\!\mathit{om}}\nolimits}
\def\SheafExt{\mathop{\mathcal{E}\!\mathit{xt}}\nolimits}
\def\Sch{\mathit{Sch}}
\def\Mor{\mathop{\mathrm{Mor}}\nolimits}
\def\Ob{\mathop{\mathrm{Ob}}\nolimits}
\def\Sh{\mathop{\mathit{Sh}}\nolimits}
\def\NL{\mathop{N\!L}\nolimits}
\def\CH{\mathop{\mathrm{CH}}\nolimits}
\def\proetale{{pro\text{-}\acute{e}tale}}
\def\etale{{\acute{e}tale}}
\def\QCoh{\mathit{QCoh}}
\def\Ker{\mathop{\mathrm{Ker}}}
\def\Im{\mathop{\mathrm{Im}}}
\def\Coker{\mathop{\mathrm{Coker}}}
\def\Coim{\mathop{\mathrm{Coim}}}
\DeclareMathSymbol{\boxtimes}{\mathbin}{AMSa}{"02}
\def\QCohstack{\mathcal{QC}\!\mathit{oh}}
\def\Cohstack{\mathcal{C}\!\mathit{oh}}
\def\Spacesstack{\mathcal{S}\!\mathit{paces}}
\def\Quotfunctor{\mathrm{Quot}}
\def\Hilbfunctor{\mathrm{Hilb}}
\def\Curvesstack{\mathcal{C}\!\mathit{urves}}
\def\Polarizedstack{\mathcal{P}\!\mathit{olarized}}
\def\Complexesstack{\mathcal{C}\!\mathit{omplexes}}
\def\Pic{\mathop{\mathrm{Pic}}\nolimits}
\def\Picardstack{\mathcal{P}\!\mathit{ic}}
\def\Picardfunctor{\mathrm{Pic}}
\def\Deformationcategory{\mathcal{D}\!\mathit{ef}}
\usepackage{cleveref}
\crefname{originaldefinition}{original definition}{original definitions}
\Crefname{originaldefinition}{Original definition}{Original definitions}
\makeatletter
\newlabel{definition-algebraic-space}{{025Y}{source}{}{sourcepointer-1}{}}
\newlabel{definition-algebraic-space@cref}{{[originaldefinition][0][]025Y}{[1][0][]source}{}{}{}}
\makeatother
\crefname{originallemma}{original lemma}{original lemmas}
\Crefname{originallemma}{Original lemma}{Original lemmas}
\makeatletter
\newlabel{lemma-base-change-representable-transformations}{{02WB}{source}{}{sourcepointer-2}{}}
\newlabel{lemma-base-change-representable-transformations@cref}{{[originallemma][0][]02WB}{[1][0][]source}{}{}{}}
\makeatother
\crefname{originallemma}{original lemma}{original lemmas}
\Crefname{originallemma}{Original lemma}{Original lemmas}
\makeatletter
\newlabel{lemma-base-change-representable-transformations-property}{{02WL}{source}{}{sourcepointer-3}{}}
\newlabel{lemma-base-change-representable-transformations-property@cref}{{[originallemma][0][]02WL}{[1][0][]source}{}{}{}}
\makeatother
\crefname{originallemma}{original lemma}{original lemmas}
\Crefname{originallemma}{Original lemma}{Original lemmas}
\makeatletter
\newlabel{lemma-composition-representable-transformations}{{02WA}{source}{}{sourcepointer-4}{}}
\newlabel{lemma-composition-representable-transformations@cref}{{[originallemma][0][]02WA}{[1][0][]source}{}{}{}}
\makeatother
\crefname{originallemma}{original lemma}{original lemmas}
\Crefname{originallemma}{Original lemma}{Original lemmas}
\makeatletter
\newlabel{lemma-representable-transformation-diagonal}{{05L9}{source}{}{sourcepointer-5}{}}
\newlabel{lemma-representable-transformation-diagonal@cref}{{[originallemma][0][]05L9}{[1][0][]source}{}{}{}}
\makeatother
\crefname{originallemma}{original lemma}{original lemmas}
\Crefname{originallemma}{Original lemma}{Original lemmas}
\makeatletter
\newlabel{lemma-representable-transformation-to-sheaf}{{02WD}{source}{}{sourcepointer-6}{}}
\newlabel{lemma-representable-transformation-to-sheaf@cref}{{[originallemma][0][]02WD}{[1][0][]source}{}{}{}}
\makeatother
\title{A functor admitting a representable morphism to an algebraic space is an algebraic space}
\author{The Stacks Project Authors}
\date{stacks-project:spaces.tex:local-snapshot-20261012; GNU FDL 1.2 or later}
\begin{document}
\maketitle
\setcounter{subsection}{1}
\noindent\textit{Editorial note.} All delivered mathematical text is the author's own, copied from the pinned Stacks Project chapter file spaces.tex; no statement or proof step was rewritten, reordered or removed.\par\medskip
\noindent\textit{Editorial note.} The statement carries the author's slogan environment; it is part of the counted statement span and is preserved verbatim, with the slogan declaration retained in the preamble.\par\medskip
\noindent\textit{Editorial note.} The six references of the counted spans are recorded as non-load-bearing navigation pointers; the argument that consumes them stays inside the counted proof span.\par\medskip
\noindent\textit{Editorial note.} Preamble: sources/stacks-spaces/retained-preamble.tex, sliced mechanically from the chapter's own preamble.tex.\par\medskip
\noindent\textit{Editorial note.} Revision r002 only changes the preamble span: the retained preamble was extended with \textbackslash{}usepackage\{verbatim\}, restored verbatim from the chapter's own preamble.tex line 8, because the counted statement uses the author's slogan environment which is defined through \textbackslash{}comment. No statement or proof text changed; r001 failed on the undefined \textbackslash{}comment.\par\medskip
\noindent\textit{Allowed standard prerequisites (editorial).}\begin{itemize}
\item A representable transformation with sheaf target has sheaf source. Lemma lemma-representable-transformation-to-sheaf (433-439), used at proof lines 1939-1940. Gives the sheaf condition for G.
\item Base change of a representable transformation is representable. Lemma lemma-base-change-representable-transformations, used at proof lines 1947-1949. Makes G x\_F G -> G x G representable.
\item The diagonal of a representable transformation is representable. Lemma lemma-representable-transformation-diagonal (456-462, produced as L44186), used at proof lines 1950-1952. Gives representability of G -> G x\_F G.
\item Composition of representable transformations. Lemma lemma-composition-representable-transformations, used at proof lines 1953-1955. Makes Delta\_G representable.
\item Base change keeps the property of being surjective and etale. Lemma lemma-base-change-representable-transformations-property, used at proof lines 1958-1960. Produces the surjective etale chart U' -> G.
\end{itemize}\medskip
\ifdefined\hypertarget\hypertarget{sourcepointer-1}{}\fi
\noindent\textit{Source reference.} Original definition 025Y: focused-seeds/AGP/sources/stacks-snapshot-20261011/tags.txt, source lines 2831--2831. Navigation pointer into the Stacks Project; the standard fact it supplies is applied inside the counted proof span and the argument resting on it stays there.\par\medskip
\ifdefined\hypertarget\hypertarget{sourcepointer-2}{}\fi
\noindent\textit{Source reference.} Original lemma 02WB: focused-seeds/AGP/sources/stacks-snapshot-20261011/tags.txt, source lines 3798--3798. Navigation pointer into the Stacks Project; the standard fact it supplies is applied inside the counted proof span and the argument resting on it stays there.\par\medskip
\ifdefined\hypertarget\hypertarget{sourcepointer-3}{}\fi
\noindent\textit{Source reference.} Original lemma 02WL: focused-seeds/AGP/sources/stacks-snapshot-20261011/tags.txt, source lines 3808--3808. Navigation pointer into the Stacks Project; the standard fact it supplies is applied inside the counted proof span and the argument resting on it stays there.\par\medskip
\ifdefined\hypertarget\hypertarget{sourcepointer-4}{}\fi
\noindent\textit{Source reference.} Original lemma 02WA: focused-seeds/AGP/sources/stacks-snapshot-20261011/tags.txt, source lines 3797--3797. Navigation pointer into the Stacks Project; the standard fact it supplies is applied inside the counted proof span and the argument resting on it stays there.\par\medskip
\ifdefined\hypertarget\hypertarget{sourcepointer-5}{}\fi
\noindent\textit{Source reference.} Original lemma 05L9: focused-seeds/AGP/sources/stacks-snapshot-20261011/tags.txt, source lines 7228--7228. Navigation pointer into the Stacks Project; the standard fact it supplies is applied inside the counted proof span and the argument resting on it stays there.\par\medskip
\ifdefined\hypertarget\hypertarget{sourcepointer-6}{}\fi
\noindent\textit{Source reference.} Original lemma 02WD: focused-seeds/AGP/sources/stacks-snapshot-20261011/tags.txt, source lines 3800--3800. Navigation pointer into the Stacks Project; the standard fact it supplies is applied inside the counted proof span and the argument resting on it stays there.\par\medskip
\begin{lemma}
\label{lemma-representable-over-space}
\begin{slogan}
A functor that admits a representable morphism to an algebraic space is
an algebraic space.
\end{slogan}
Let $S$ be a scheme contained in $\Sch_{fppf}$.
Let $F$ be an algebraic space over $S$.
Let $G \to F$ be a representable transformation of functors.
Then $G$ is an algebraic space.
\end{lemma}

\begin{proof}
By Lemma \ref{lemma-representable-transformation-to-sheaf}
we see that $G$ is a sheaf. The diagram
$$
\xymatrix{
G \times_F G \ar[r] \ar[d] & F \ar[d]^{\Delta_F} \\
G \times G \ar[r] & F \times F
}
$$
is cartesian. Hence we see that $G \times_F G \to G \times G$
is representable by
Lemma \ref{lemma-base-change-representable-transformations}.
By
Lemma \ref{lemma-representable-transformation-diagonal}
we see that $G \to G \times_F G$ is representable.
Hence $\Delta_G : G \to G \times G$ is representable as a composition
of representable transformations, see
Lemma \ref{lemma-composition-representable-transformations}.
Finally, let $U$ be an object of $(\Sch/S)_{fppf}$
and let $U \to F$ be surjective and \'etale. By assumption
$U \times_F G$ is representable by a scheme $U'$. By
Lemma \ref{lemma-base-change-representable-transformations-property}
the morphism $U' \to G$ is surjective and \'etale. This verifies
the final condition of Definition \ref{definition-algebraic-space} and we win.
\end{proof}

\end{document}
